



\section{Experiments}
\label{sec:exp}
In this section, we present a comprehensive evaluation of \modelname{}. Evaluation method are described in Section~.\ref{sec:eval_metrics}, Qualitative results are provided in Section~\ref{sec:Qualitative_Results}, quantitative analyses in Section~\ref{sec:Quantitative_Results}, and ablation studies examining the effect of individual design choices are discussed within the Method Section~\ref{sec:method}.



\subsection{Evaluation metrics.}
\label{sec:eval_metrics}
We evaluate \modelname{} using four complementary protocols: PRISM-Bench-licensed~\cite{fang2025flux}, GenEval~\cite{ghosh2023geneval}, TaBR (Section~\ref{sec:eval}), and overall user preference—covering prompt adherence, compositional reasoning, aesthetics, and expressiveness. 
PRISM-Bench-licensed is an automatically derived variant of PRISM-Bench in which prompts referencing copyrighted or unlicensed concepts are filtered out, ensuring a fair comparison as \modelname{} is trained on 100\% licensed data (see Section~\ref{sec:ethics}). PRISM-Bench was designed to expose alignment gaps that saturate on earlier benchmarks; for example, GenEval’s single-object category often reaches near-perfect scores, masking differences across models. 
For completeness, we still report GenEval, which probes object count, color attribution, and spatial relations with controlled prompts. We follow the official protocols: PRISM is scored automatically with GPT-4.1-Vision, and GenEval using detectors and CLIP-based similarity. 


TaBR measures long-caption expressiveness via a caption$\rightarrow$generation$\rightarrow$reconstruction loop with human raters comparing reconstructions to the original image (Section ~\ref{sec:eval}). Annotators are shown the original and two reconstructions (from different models) and asked: \textit{``Which Image is more similar to the one above?''}  We perform this measurement on a test-set of $60$ image that are not part of our training data. We also report overall user preference using blinded, side-by-side comparisons, asking: \textit{``Which image do you prefer?''} The evaluation covers a test set of 150 captions spanning diverse visual and stylistic domains. For both studies, image order and model identity are randomized per trial. We aggregate responses by computing, for each model, the percentage of questions won (win rate) and the percentage of ties (draw rate). In total, we collect over 5{,}000 votes to ensure a robust evaluation.
%A visual example of the user study is provided in the appendix.



% \eliran{Overall user preference directly measures human perceptual quality and alignment. Annotators are presented with image pairs generated by \modelname{} and a baseline model using the same caption and are asked to select which image they prefer overall. Each comparison is performed anonymously to avoid bias. The evaluation covers a test set of 150 captions spanning diverse visual and stylistic domains. We report the mean user preference for \modelname{} over each baseline.}\ekc{If this is standard then we can remove it}
% \ekc{also, both in TaBR and user preference we should mention the number of users questioned per pair}

\subsection{Qualitative Results}
\label{sec:Qualitative_Results}


Figure~\ref{fig:tabr} shows qualitative examples from the TaBR evaluation set, with original reference images alongside reconstructions produced by FLUX, Qwen-Image, and \modelname{}. As illustrated, \modelname{} achieves the closest resemblance to the originals, preserving global structure and fine-grained details. In particular, it maintains pose, camera angle, and overall scene layout more faithfully than the baselines.

Figure~\ref{fig:contextual} presents \textit{Contextual Contradictions} drawn from ContraBench~\cite{huberman2025image} and Whoops~\cite{Bitton_Guetta_2023}, which probe robustness to atypical relations, whereas prior models often default to plausible but incorrect configurations. \modelname{} preserves the full set of specified attributes and relations, e.g., bear performing handstand and women writing with a dart. We attribute this fidelity to training on structured captions that explicitly encode concepts and context. Additional visual results are provided in the Appendix, Figure~\ref{fig:app_samples}, highlighting \modelname{}’s creative range. Figure~\ref{fig:app_disentanglement} provide additional examples for controllable refinement.







%In most cases, Qwen-Image appears to satisfy the textual prompts at first glance; however, closer inspection reveals inconsistencies. In the balloon example, the package is shown with a balloon above it, yet the loose string indicates that the balloon does not actually lift the package. 
%In the vase example, the “upside down” description is fulfilled by flipping the entire scene—vase and table included—rather than inverting only the bouquet as intended.

\subsection{Quantitative Results}
\label{sec:Quantitative_Results}

Table~\ref{tab:prism} reports results on PRISM-Bench-licensed. \modelname{} attains the highest alignment score among open-source models, substantially narrowing the prompt-alignment gap to closed systems such as Gemini~2.5~Flash~Image. Despite operating at a moderate 8B parameter scale, \modelname{} rivals larger models in prompt understanding and text–image grounding, underscoring the impact of structured-caption supervision.
We additionally report scores for \modelname{} when paired with a closed-source VLM~\cite{comanici2025gemini} and with our fine-tuned open-source alternative. Surprisingly, the open-source variant achieves competitive results while relying on a significantly smaller and more cost-effective VLM to produce the structured prompts.

On GenEval (Table~\ref{tab:geneval}), \modelname{} delivers consistently strong results across tasks, with a pronounced gain in positional understanding, a known weakness for most models~\cite{ghosh2023geneval}. We attribute this improvement to training on structurally detailed captions that explicitly encode spatial relations.

Table~\ref{tab:tabr-comparison} presents TaBR results. \modelname{} achieves the highest reconstruction fidelity across baselines, winning over 90\% of comparisons against SD3.5-Large and 88.9\% against HiDream-II-Full, while maintaining strong performance against FLUX and Qwen-Image. These findings indicate superior expressive power derived from long structured-caption training. Finally, in overall user preference (Table~\ref{tab:user-preference}), \modelname{} is consistently favored over SD3.5, FLUX.1-dev, and FLUX.1-Krea-dev, and attains comparable ratings to Qwen-Image and HiDream-II-Full—demonstrating competitive perceptual quality alongside strong prompt alignment.






\begin{table}
{\footnotesize
\setlength{\tabcolsep}{5pt}
  \centering
  \begin{tabular}{lccccc}
    \toprule
    & \textbf{SD3.5} & \textbf{FLUX-dev} & \textbf{FLUX-Krea} & \textbf{HiDream} & \textbf{Qwen} \\    
    \midrule
    \modelname{} & \textbf{90.5\%} & \textbf{76.9\%} & \textbf{66.4\%} & \textbf{88.9\%} & \textbf{59.2\%} \\
    Baseline & 2.9\% & 8.0\% & 13.9\% & 4.5\% & 22.0\% \\
    Draw & 6.6\% & 15.1\% & 19.7\% & 6.6\% & 18.8\% \\
    \bottomrule
  \end{tabular}
  }
\caption{\textbf{Text-as-a-Bottleneck Reconstruction (TaBR) evaluation.} Win rates of \modelname{} against each model. \modelname{} outperforms all open-source baselines, reflecting greater expressive power due to training on long structured captions.}

  \label{tab:tabr-comparison}

\end{table}

%\vspace{-0.4cm}


% \begin{table}
% {\footnotesize
% \setlength{\tabcolsep}{6pt}
%   \centering
%   \begin{tabular}{lccccc}
%     \toprule
%     & \textbf{SD3.5-Large} & \textbf{FLUX.1-dev} & \textbf{FLUX.1-Krea-dev} & \textbf{HiDream-I1-Full} & \textbf{Qwen-Image} \\    
%     \midrule
%     FIBO & \textbf{90.5\%} & \textbf{76.9\%} & \textbf{66.4\%} & \textbf{88.9\%} & \textbf{59.2\%} \\
%     Baseline & 2.9\% & 8.0\% & 13.9\% & 4.5\% & 22.0\% \\
%     Draw & 6.6\% & 15.1\% & 19.7\% & 6.6\% & 18.8\% \\
%     \bottomrule
%   \end{tabular}
%   }
%     \caption{User study results for the Text-as-a-Bottleneck Reconstruction (TaBR) evaluation.
% Win rates of FIBO over each baseline method.}
%   \label{tab:tabr-comparison}

% \end{table}
\begin{table}
{\footnotesize
\setlength{\tabcolsep}{4.5pt}
  \centering
  \begin{tabular}{lccccc}
    \toprule
     & \textbf{SD3.5} &  \textbf{FLUX-dev} & \textbf{FLUX.1-Krea} & \textbf{HiDream} & \textbf{Qwen} \\    
    \midrule
    \modelname{} & \textbf{69.1\%} & \textbf{58.6\%} & \textbf{47.5\%} & 42.3\% & \textbf{42.1\%} \\
    Baseline & 19\% & 30.4\% & 35.1\% & \textbf{43.5\%} & 40.2\% \\
    Draw & 11.9\% & 11\% & 17.4\% & 19.7\% & 17.7\% \\
    \bottomrule
  \end{tabular}
  }
  \caption{\textbf{Overall user preference.} \modelname{} clearly surpasses SD3.5, FLUX.1-dev and Flux.1-Krea-dev, and achieves comparable performance to Qwen-Image and HiDream.}
  \label{tab:user-preference}
\end{table}